\section{Ablation Studies}
\label{sec:ablation}

To identify which components of SheafDB contribute most to end-to-end performance, we conduct ablation studies using the optimization framework introduced in Section~\ref{sec:optimization}.

\paragraph{Optimization framework.}
The framework defines 33 independent optimization modules distributed across three layers: sheaf layer (13 modules covering interval indexing, cocycle caching, morphism inlining, restriction-map memoization, section merging strategies), KG-compatibility layer (10 modules governing reification bypass, triple---tuple bridge strategies, join planning), and query layer (10 modules including predicate pushdown, result pipelining, aggregation fusion, and index-only scans).

\paragraph{Configuration profiles.}
Six pre-defined profiles span the optimization space:
\begin{itemize}
  \item \textbf{Minimal (7 on)} --- core tuple storage only; no caching, no indexing beyond primary key.
  \item \textbf{Default (22 on)} --- production configuration with interval indexing, cocycle caching, and predicate pushdown enabled.
  \item \textbf{Research (27 on)} --- default plus experimental modules: speculative restriction, morphism inlining, and lazy section merge.
  \item \textbf{Max (33 on)} --- all modules enabled, including debugging instrumentation and profiling hooks.
  \item \textbf{Memory (10 on)} --- optimized for peak memory efficiency: indexing disabled, streaming merge only.
  \item \textbf{Debug (0 on)} --- all modules disabled; raw tuple storage with safety checks.
\end{itemize}

\paragraph{Ablation strategies.}
Three complementary strategies are employed:
\begin{enumerate}
  \item \textbf{Leave-one-out.} Each of the 33 modules is disabled individually while all others remain at the default profile, measuring the performance delta on each workload category.
  \item \textbf{Pairwise.} All $\binom{33}{2} = 528$ pairs are independently toggled to detect interaction effects (synergistic or antagonistic module combinations).
  \item \textbf{Random-subset.} 1\,000 random configuration samples across the full 33-dimensional binary space provide a global sensitivity landscape.
\end{enumerate}

\paragraph{Sensitivity sweeps.}
Ten continuous parameters---including interval-tree node size, cocycle cache TTL, restriction-map batch threshold, and section-merge fan-in---are swept across two orders of magnitude each, producing per-parameter Pareto fronts.

\paragraph{Scalability dimensions.}
Nine dimensions---tuple count, arity, context overlap ratio, provenance chain depth, temporal window count, attribute cardinality, number of sources, query concurrency, and batch size---are varied independently to characterize robustness.

\paragraph{Expected insights.}
Preliminary analysis indicates that interval indexing (C5--C6), cocycle caching (C8), and reification bypass (C3--C4) are the three highest-impact modules, together accounting for more than 60\% of the performance gap between Minimal and Default profiles. Pairwise sweeps reveal a synergistic effect between morphism inlining and restriction-map memoization that yields super-linear speedups on provenance workloads (C7).

\begin{figure}
  \centering
  \includegraphics[width=0.48\textwidth]{figures/ablation_heatmap.pdf}
  \caption{Leave-one-out ablation heatmap. Rows are workloads; columns are optimization modules. Color encodes speedup factor relative to the Minimal profile.}
  \label{fig:ablation_heatmap}
\end{figure}

\begin{figure}
  \centering
  \includegraphics[width=0.48\textwidth]{figures/sensitivity_radar.pdf}
  \caption{Sensitivity radar across 9 scalability dimensions for the Default configuration. Each axis is log-scaled; values are normalized to the peak.}
  \label{fig:sensitivity_radar}
\end{figure}

\input{tables/ablation_summary.tex}
